\subsection{Extracting an actual zero name from a positive count}
\label{sec:zero-extraction}

A positive contour count contains more constructive information than an
unrelated sequence of small residuals. Here is a coherent root extraction,
using only the quantitative minimum-modulus construction and the finite
winding rules already proved.

Fix one rational rectangle $Q_0$ with a certified zero-free boundary and
positive xi winding. Choose an integer $A\ge1$ so that
$|s-2|\le A$ throughout $Q_0$, and use $f(z)=\xi(2+Az)$. The entire
theta-integral construction gives, with $R=4A+2$,
\[
        |f(z)|\le M:=1+2R(R+1)R!\quad(|z|\le4),
             \qquad f(0)=\xi(2)=\zeta(2)/\pi\ge1/4.
\]
The same expression with $R=5A+2$ bounds the larger disk used for Taylor
approximation. The minimum-modulus lemma therefore applies with the fixed
integer $K=4\lceil\log_2(16M)\rceil$. In its notation one may take
\[
       L_\delta=2^{-83}M^{-20}(\delta/6)^K,
\]
so that outside at most $K$ disks of radius $\delta$ the normalized
function has modulus at least $3L_\delta/4$. The global derivative on
$Q_0$ is bounded by a fixed primitive recursive integer $D\ge1$; for
example the Cauchy bound from the displayed disk estimate permits $D=M$.
These constants may be very large, but none is obtained by an unbounded
search.

For a current rectangle $[a,a+w]\times[b,b+h]$, use the $K+1$ candidates
\[
 x_j=a+w/3+\frac{(j+1)w}{3(K+2)},\qquad
 y_j=b+h/3+\frac{(j+1)h}{3(K+2)}\quad(0\le j\le K),
\]
and put $\delta=\min(w,h)/(9A(K+2))$. Twice this radius is smaller
than either candidate spacing in normalized coordinates, so each exceptional
disk excludes at most one vertical candidate and at most one horizontal
candidate. There is therefore a safe full segment in each direction.
On every candidate segment evaluate $\xi$ to error $L_\delta/16$ on a
mesh of physical width at most $L_\delta/(8D)$. Choose the first segment
whose rational samples have modulus at least $L_\delta/2$, tested by
squaring. A safe candidate passes this test, and every passing candidate
has true modulus at least $5L_\delta/16$ on the entire segment. Record
the smaller margin $L_\delta/4$ for each of the two chosen cuts.

If the parent's recorded margin is $m$, give every child the common margin
\[
                         m'=\tfrac12\min(m,L_\delta/4).
\]
Every child edge then has true
modulus at least $2m'$. Restrict the inherited boundary certificates
to the child edges, and refine both inherited edges and the two cuts to
a common mesh of width at most $m'/(8D)$, with sample error at most $m'/8$.
The resulting Lipschitz certificates verify the claimed margin $m'$.
Reuse
identical samples with opposite orientations on shared child edges. This
constructs full closed-boundary certificates for all four children, not
just certificates for their new cuts. Each side length lies between
$1/3$ and $2/3$ of the corresponding parent side length.

Each child winding is nonnegative. To see this without already assuming
analytic root existence, approximate $\xi$ on that child's boundary by a
rational complex polynomial to less than half the certified boundary
margin. The finite zero-free homotopy identifies the two windings. The
polynomial engine of Section~\ref{sec:polynomial-algebra} factors that
polynomial, so its winding is the sum of the nonnegative indices of its
linear factors. Cancellation of the new internal edges makes the four
child windings sum to the parent's winding. At least one is consequently
positive; choose the first such child in a fixed ordering.

Iterating this procedure gives nested rational rectangles $Q_j$ with
positive winding and diameter at most $D_0(2/3)^j$, for a supplied
$D_0\ge1$. All searches are bounded, and the bounds are recomputed by
primitive recursion from the current rational rectangle and its inherited
boundary margin. In particular this is not a search that waits until a
possibly vanishing sample becomes visibly nonzero.

Let $c_j$ be the centre of $Q_j$. Then
\[
                        |\xi(c_j)|\le D\,\operatorname{diam}(Q_j).
\]
Indeed, a strict violation would put the entire boundary image in a disk
centred at $\xi(c_j)$ which misses zero; its winding would be zero. The
finite-precision version proves the displayed weak inequality by ruling
out every rationally witnessed violation, using the same derivative bound.

Choose an integer $L$ with $2^L\ge4DD_0$, and set
\[
                         \alpha(k)=c_{\,2(k+L)}.
\]
Since $(2/3)^2<1/2$, nestedness gives, for $\ell\ge k$,
\[
       |\alpha(\ell)-\alpha(k)|\le2^{-k-2},\qquad
       |\xi(\alpha(k))|\le2^{-k-2}.
\]
The centres therefore form one primitive recursive fast Cauchy name for an
actual zero. The initial positive boundary margin and derivative bound
also keep this zero positively separated from $\partial Q_0$, so it lies
inside the original rectangle. Every recurrence invariant used to obtain these inequalities
concerns finite rational rectangles, contour certificates, signed integer
counts and fixed-precision approximation errors. No countable choice or
decision of exact zero ordinates is involved.

For rectangles above the real axis, the explicitly nonvanishing factors
relating $\xi$ and $\zeta$ identify the extracted point as a nontrivial
zeta zero. Applying the construction to disjoint certified height bands
gives pairwise apart zero names. This proves constructive infinitude when
the counting estimate makes those band counts positive; it does not supply
a procedure for deciding equal ordinates or sorting all exact zeros.
